\chapter{上下文学习}
\label{lab:91_prompt_reasoning}

\noindent\textbf{配套文件：}\path{notebook/91_prompt_reasoning.ipynb}\par
\noindent\textbf{教材关联：}\bookxrefshort{ch:context-learning}。

\section*{实验目标}
控制任务、模型和预算后比较提示策略。

\section*{准备及定位}
先阅读本册实验总则，确认Notebook中的依赖与资产单元。原理计算、库接口迁移和真实模型训练可能具有不同资源要求，应分别记录设备、版本及运行范围。

源码定位可从下列函数开始；应结合前置数据与配置单元阅读。
\begin{itemize}
\item \texttt{my\_terms}
\item \texttt{my\_example\_score}
\item \texttt{my\_select\_examples}
\end{itemize}

\section*{操作及观察}
\begin{enumerate}
\item 冻结问题集、输出格式及最终答案抽取规则。
\item 对照示例选择、直接回答与推理提示，保存每次提示的完整身份。
\item 在相同预算口径下进行多样本生成与Self-Consistency聚合。
\item 比较正确率、格式失败及成本，检查错误抽取和无效样本的分母处理。
\end{enumerate}

\section*{结果判读及提交}
提交逐题结果、样本数、提示版本和聚合规则。多数一致不保证事实正确；长推理文本也不构成正确性的独立证明。

提交时附上所执行单元范围、环境与制品身份、原始结果及失败说明。将未运行项目明确列出，不能用Notebook中已有输出替代本次执行证据。
